โครงงานนี้นำเสนอการประยุกต์ใช้กระบวนการวิทยาการข้อมูลในการวิเคราะห์ปัจจัยที่มีผลต่อผลการแข่งขันฟุตบอลพรีเมียร์ลีกอังกฤษ และการพัฒนาแบบจำลองการเรียนรู้ของเครื่องสำหรับทำนายผลการแข่งขันในรูปแบบ 3 ผลลัพธ์ ได้แก่ ทีมเจ้าบ้านชนะ เสมอ และทีมเยือนชนะ

การดำเนินงานได้บูรณาการชุดข้อมูลการแข่งขันฟุตบอลย้อนหลัง ตั้งแต่ฤดูกาล 2014 ถึง 2026 จากหลากหลายแหล่งข้อมูล ครอบคลุมสถิติผลการแข่งขัน รายชื่อผู้เล่นที่ลงสนาม ประวัติตัวผู้เล่นและการบาดเจ็บ รวมถึงค่าความคาดหวังในการทำประตู (Expected Goals: xG) ข้อมูลทั้งหมดผ่านกระบวนการทำความสะอาด การเชื่อมโยงโครงสร้างข้อมูล และการวิศวกรรมฟีเจอร์เพื่อสร้างตัวแปรเชิงทำนายสำคัญ ได้แก่ คะแนนความแข็งแกร่ง Elo Rating, ค่า xG, ฟอร์มการเล่นย้อนหลัง, อัตราต่อรองจากตลาดเดิมพัน และดัชนีการขาดหายไปของผู้เล่นหลัก จากนั้นทำการฝึกฝนและประเมินประสิทธิภาพของแบบจำลองการจำแนกประเภท ได้แก่ Logistic Regression, Decision Tree และ Random Forest

ผลการศึกษาพบว่า การนำข้อมูลระดับผู้เล่นร่วมกับตัวชี้วัดเชิงสถิติขั้นสูง (xG) และคะแนนความแข็งแกร่งของทีมมาบูรณาการ ช่วยเพิ่มประสิทธิภาพของแบบจำลองในการทำนายผลการแข่งขันได้อย่างมีนัยสำคัญ เมื่อเปรียบเทียบกับการใช้เพียงสถิติการแข่งขันย้อนหลังแบบดั้งเดิม โดยแบบจำลอง Random Forest ให้ผลการทำนายที่มีประสิทธิภาพสูงที่สุด และมีความสอดคล้องกับแนวโน้มของตลาดอัตราต่อรอง ตลอดจนสามารถระบุระดับความสำคัญของฟีเจอร์ที่มีผลกระทบต่อผลการแข่งขันได้อย่างชัดเจน องค์ความรู้และกรอบความคิดจากโครงงานนี้สามารถนำไปประยุกต์ใช้เป็นแนวทางในการประมวลผลข้อมูลกีฬาที่มีความซับซ้อนหลายมิติ และสนับสนุนการตัดสินใจเชิงวิเคราะห์ในอุตสาหกรรมกีฬาต่อไป